%======================================================================
%   PRESENTATION OF CAT SCHEMES --- STATE OF THE ART
%
\section{Compact Approximate Taylor (CAT) schemes} \label{sec:CAT} 
%\rred{The notation used in this section is the one corresponding to scalar problems: it should be changed to capital bold letters or to add a sentence to say that the description is done in the scalar case to simplify the notation...} \ema{Done}

The focus of this chapter is the presentation of a family of numerical methods for non-linear systems of conservation law, named Compact Approximate Tayor (CAT) schemes\footnote{For simplicity we will summarize the numerical deduction of the method in the scalar case.}.
This family is based on an approximate Taylor procedure that constitutes a proper generalization of Lax-Wendroff (LW) method, in the sense that it reduces to the standard high-order LW method when the flux is linear. 
%
In this section we recall the LW and CAT procedures of second and fourth order.

\subsection{Lax Wendroff procedure} \label{ssec:LW}
A scheme of historical as well as practical importance is the celebrated \textbf{Lax Wendroff} scheme introduced by Peter Lax and Burton Wendroff in 1960 in \cite{LaxWendroff}, and \cite{Lax1957,Wendroff,LaxWendroffHigh,Hirsch22007,Toro2009}. It has been the most widely adopted scheme for aeronautical applications, up to the end of the 1980s under various forms.\footnote{Probably the most widely used variant is the two-step Mac Cormack scheme, which has properties similar to LW schemes, and avoids the computation of the second derivative. 
Published originally at a conference in 1969, the paper has been reproduced in \cite{maccormack2002effect}.}
%\giovanni{[I think we should cite the Mac Cormack scheme, which is based on a predictor-corrector formulation, gives results comparable to Lax-Wendroff, and avoids the computation of the second derivative. It is NOT Runge-Kutta type, and is probably more used by aerospace engineers that LW. I try to cite it as a footnote. By the way, given the efficiency of MC scheme with respect to LW, it would be interesting to explore its high order generalization...]}



%
The original derivation of Lax and Wendroff was based on a Taylor expansion in time of function $u$ at point $(x_i,t)$ up to second order of accuracy, thus
%so as to achieve second order accuracy in time. Thus,
\begin{equation}
    \label{sec:LW_Taylor2}
    u(x_i,t+\Delta t)=u(x_i,t) + \Delta t \, u_t(x_i,t) + \frac{\Delta t^2}{2}u_{tt}(x_i,t) + O(\Delta t^3),
\end{equation}
where $\Delta t>0$ is a small increment in time.
The numerical scheme is then obtained by neglecting the higher order term in $\Delta t$, using the governing equation to replace time derivatives by spatial ones, and then substituting the obtained space derivatives with their finite difference approximations. 
%
For the \textit{linear case}, $u_t =-a u_x$ with $f(u)=au$, we obtain
\bea
    \label{eq:lw0}
    u(x_i,t^n+\Delta t)=u(x_i,t^n) - \Delta t \, a u_x(x_i,t^n)  + \frac{\Delta t^2}{2} a^2 u_{xx} (x_i,t) + O(\Delta t^3).
\eea
Using centred finite differences to approximate spatial derivatives, the numerical scheme follows as:
\begin{equation}
    \label{lw}
    u_i^{n+1}=u_i^n - \Delta t \, a \, \frac{u_{i+1}^n-u_{i-1}^n}{2\Delta x} + \frac{\Delta t^2}{2} a^2  \Frac{u_{i+1}^n - 2u_i^n + u_{i-1}^n}{\Delta x^2},
\end{equation}
where 
%$\{x_i\}$ are the nodes of a uniform mesh of step $\Delta x;$ 
$u_i^m$ is an approximation of the point value of the solution at position $x_i$ at the time $t^m$. %$n\Delta t$, in which $\Delta t$ is the time step. 
A useful alternative formulation written in conservative form yields
\begin{equation}
    \label{lw_conservative}
    u_i^{n+1}=u_i^n - \frac{\Delta t}{\Delta x}\left( F_{i+\ha}^{\rm{LW}}- F_{i-\ha}^{\rm{LW}} \right),
\end{equation}
where the so-called LW numerical flux, $F_{i+\ha}^{\rm{LW}}$, is given by
\bea \label{eq:LW_flux}  
F_{i+\ha}^{\rm{LW}} = \frac{a}{2}\left(u_{i+1}^n + u_i^n\right) - \frac{a^2\Delta t}{2\Delta x}\left(u_{i+1}^n - u_i^n\right).
\eea
%which emphasizes the conservative structure of the method.
%
%\ema{
The \textit{non-linear case}, $u_t =-f_x(u)$, yields
\begin{equation}
    \label{non-lin_lw}
    u_i^{n+1}=u_i^n - \frac{\Delta t}{2\Delta x}\left(f_{i+1}^n-f_{i-1}^n\right) + \frac{\Delta t^2}{2\Delta x^2}\left(\text{A}_{i+\ha}\left(f_{i+1}^n - f_{i}^n\right) - \text{A}_{i-\ha}\left(f_{i}^n - f_{i-1}^n\right)\right),
\end{equation}
where $f_{i+j}^n = f(u_{i+j}^n)$ for $j = -1,0,1$ and $\text{A}$ is an approximation of the derivative of $f$, i.e $A= \partial f/\partial u$. 
Hence, $\text{A}_{i\pm\ha}$ is the approximation derivative of $f$
evaluated at $u_{i+\ha}^n=\ha(u_{i}^n + u_{i\pm1}^n)$, that is $\text{A}_{i\pm\ha}\equiv\text{A}(u_{i+\ha}^n)$, or, the average between the cell-based derivative, that is $\text{A}_{i\pm\ha}\equiv \ha \left( \text{A}(u_i^n)+\text{A}(u_{i\pm1}^n) \right)$. 
%}
%\rred{There is again the same problem with the notation: if the numerical method is described for systems,  bold capital letters should be used for everything and not only for the Jacobian. If, instead, the numerical methods is described for scalar problems, then $A_{i + 1/2}$ is an approximation of the derivative of $f$...} \ema{Done}
Notice that they  depend non-linearly on variable $u$. Moreover they are always evaluated at time $t^n$, so we can omit this time dependency.
The alternative conservative formulation of the LW scheme is expressed as:
\begin{equation}
    \label{non-lin_lw_conservative}
    u_i^{n+1}=u_i^n - \frac{\Delta t}{\Delta x}\left(F_{i+\ha}^{\rm{LW}}- F_{i-\ha}^{\rm{LW}}\right),
\end{equation}
where the numerical flux
\bea \label{eq:flux_NL} 
    F_{i+\ha}^{\rm{LW}} = \underbrace{\frac{1}{2}\left( f_{i+1}^n + f_i^n \right)}_{\text{Physical flux}} \underbrace{- \frac{\Delta t}{2\Delta x}A_{i+\ha}\left(f_{i+1}^n - f_i^n\right)}_{\text{Dissipation}} ,
\eea
is composed of two parts: the average of the physical fluxes at cells $i$ and $i+1$, and, the numerical dissipation.





\subsection{Compact Approximate Taylor (CAT) procedure} \label{ssec:CAT}
The generalized Lax-Wendroff method is used to update the numerical solution:
\bea \label{eq:general_LW}
 u_i^{n+1} = u_i^n + \sum_{k=1}^{2P} \frac{(\Delta t)^k}{k!} u_i^{(k)},
\eea
%where $\{x_i\}$ are the nodes of a uniform mesh of step $\Delta x;$ 
where we recall that
$u_i^n$ is an approximation of the value of the exact solution $u(x,t)$ at time $t^n$ at position $x_i$ \cite{Hirsch12007}, and, $u_i^{(k)}$ is an approximation of $\partial_t^k u(x_i,t^n)$. 
The $k$-th derivative in time of $u$ is computed with a compact and numerical version of the Cauchy-Kovalesky procedure introduced by Carrillo and Par{\'e}s in \cite{Carrillo-Pares}.  \\ 
The final expression of the $2P$-order CAT method in conservative form is:
\begin{equation}
    \label{CAT2P}
    u_i^{n+1} = u_i^n + \frac{\Delta t}{\Delta x}\left(F_{i-\ha}^P - F_{i+\ha}^P\right).
\end{equation} 
Let us introduce the sets $\mathcal{S}_{i\pm\ha}^P$ of values $u_i^n$ on  
stencils centered around interface $i\pm\ha$ of size $2P$, 
that is
\bea \label{eq:stencil2P} 
%\mathcal{S}_{i-\ha}^P = \left\{ u_{i-P}^n,\ldots, u_{i-1}^n, u_{i}^n, \ldots u_{i+P-1}^n \right\}, \quad
\mathcal{S}_{i+\ha}^P = \left\{ u_{i-P+1}^n,\ldots, u_{i}^n, u_{i+1}^n, \ldots u_{i+P}^n \right\}.
\eea
%\rred{There is no need to introduce the stencil for both $i \pm 1/2$: $i + 1/2$ is enough...}
The flux functions $F_{i\pm \ha}^P$ are then computed, respectively, on the sets $\mathcal{S}_{i\pm\ha}^P$, as 
\begin{equation}
    \label{FP}
    F_{i+\ha}^P = \sum_{k=1}^{2P}\frac{\Delta t^{k-1}}{k!}f^{(k-1)}_{i+\ha}, %=  \sum_{k=1}^{2P}\frac{\Delta t^{k-1}}{k!} \athcal{A}_{P}^{0,\ha}\left(f^{(k-1)}_{i,*},\Delta x\right)
\end{equation}
and 
\bea \label{eq:interpolA} 
f^{(k-1)}_{i+\ha} = \mathcal{A}_{P}^{0,\ha}\left(f^{(k-1)}_{i,*},\Delta x\right), 
\quad 
\text{with}
\quad
\mathcal{A}_{P}^{0,\ha}\left(f^{(k-1)}_{i,*},\Delta x\right)  = \sum_{p=-P+1}^P\gamma_{P,p}^{0,\ha} \, f_{i+p}^{(k-1)},
\eea 
where $\mathcal{A}_{P}^{0,\ha}$ is an interpolation formula of order $2P-1$ based on $2P$-point stencil.
In the following we use the index $i$ for the cell global index, $j$ for the local position inside the stencil, $r$ for the Taylor expansion in time, and, $k$/$(k)$ to refer to the $k$-th time step/$(k)$th time derivative. 
%\rred{The notation for interpolation, numerical differentiation, and the use of the $*$ symbol should be introduced before in section 2.4...} \ema{Done}
For the sake of clarity, we detail the description of the second order ($P=1$) CAT2, in the next sub-section, and CAT4 in the Appendix \ref{sec:CAT4}. 


\subsubsection{Second order version -- CAT2}
Let $p\in \mathbb{N}$ denote an integer such that $0 \leq p \leq P$. 
In the case $P=1$ then the relative stencil is simply $\mathcal{S}^1_{i+\ha} = \{u_i^n,u_{i+1}^n\}$ for interface $i+1/2$,
while the flux reconstructions are:
\bea \label{eq:fluxCAT2}
	F_{i+\ha}^1 = f^{(0)}_{i+\ha} + \frac{\Delta t}{2}f^{(1)}_{i+\ha}, \qquad
	F_{i-\ha}^1 = f^{(0)}_{i-\ha} + \frac{\Delta t}{2}f^{(1)}_{i-\ha},
\eea
where $f^{(0)}_{i+\ha} = \ha\left(f_i^n + f_{i+1}^n\right)$ 
%and $f^{(0)}_{i-\ha} = \ha\left(f_i^n + f_{i-1}^n\right)$ 
is the interpolation of the flux at time $t^n$, while
$f^{(1)}_{i+\ha}$  %= \ha\left(f_i^{(1)} + f_{i+1}^{(1)}\right)$ 
%and $f^{(1)}_{i-\ha}$ %= \ha\left(f_i^{(1)} + f_{i-1}^{(1)}\right)$ 
is the interpolation of the first time-derivative of the flux for any interface $i+1/2$.
This implies that $\gamma_{1,p}^{0,\ha}=\ha$ in (\ref{eq:interpolA}).
These are computed over stencils $\mathcal{S}_{i\pm\ha}^1$ for
$p\in \{ 0,1 \}$ as
\begin{align}
\nonumber
\mathcal{S}_{i-\ha}^1 &:&
f^{(1)}_{i-\ha} = \frac{f_{i-1}^{(1)}+f_i^{(1)} }{2}, \quad 
&f_{i-1,p}^{(1)} = \frac{f\left( u_{i-1+p}^n + \Delta t \, u_{i-1,p}^{(1)}\right) - f_{i-1+p}^n}{\Delta t}, 
 & u_{i-1,p}^{(1)} = -\frac{f_{i}^n - f_{i-1}^n}{\Delta x},\\
\nonumber
\mathcal{S}_{i+\ha}^1 &:& 
f^{(1)}_{i+\ha} = \frac{f_i^{(1)} + f_{i+1}^{(1)}}{2}, \quad 
&f_{i,p}^{(1)} = \frac{f\left( u_{i+p}^n + \Delta t \, u_{i,p}^{(1)}\right) - f_{i+p}^n}{\Delta t}, 
 &  u_{i,p}^{(1)} = -\frac{f_{i+1}^n - f_{i}^n}{\Delta x}.
\end{align}
Notice that the $\mathcal{S}_{i+\ha}^1$ expressions can be obtained from $\mathcal{S}_{i-\ha}^1$ by replacing $i$ by $i+1$.
%\giovanni{[There is no need to specify two expressions for $f^{(1)}_{i\pm \frac12}$, since the second one is obtained from the first one by  replacing $i$ by $i+1$. Same thing can be said for Eqs.\eqref{F2+} and \eqref{F2-}. Alternatively, we may leave both expressions, and emphasize that one can be obtained from the other by replacing, say, $i$ by $i+1$ (translation invariance).]}
\begin{comment} 
\begin{itemize}
    \item The relative stencils are:
            $$\mathcal{S}^1_{i+\ha} = \{u_i^n,u_{i+1}^n\} \quad {\rm{and}}         \quad \mathcal{S}^1_{i-\ha} = \{u_{i-1}^n,u_i^n\};$$ 
    \item The flux reconstructions are:
        \begin{align}
            F_{i+\ha}^1 &= f^{(0)}_{i+\ha} + \frac{\Delta t}{2}f^{(1)}_{i+\ha}, \\
            F_{i-\ha}^1 &= f^{(0)}_{i-\ha} + \frac{\Delta t}{2}f^{(1)}_{i-\ha};
        \end{align}
        where 
        \begin{enumerate}
            \item $f^{(0)}_{i+\ha} = \ha\left(f_i^n + f_{i+1}^n\right)$ and $f^{(0)}_{i-\ha} = \ha\left(f_i^n + f_{i-1}^n\right)$ are the interpolations of the stage values at time $t^n;$  
            \item $f^{(1)}_{i+\ha} = \ha\left(f_i^{(1)} + f_{i+1}^{(1)}\right)$ and $f^{(1)}_{i-\ha} = \ha\left(f_i^{(1)} + f_{i-1}^{(1)}\right)$ are the interpolations of the first time derivative of the flux at time $t^n;$
            \begin{itemize}
                \item[a1.] $f_{i-1+j}^{(1)} = \frac{1}{\Delta t} \left(f\left(U_{i-1+j}^n + \Delta tU_{i-1,j}^{(1)}\right) - f_{i-1+j}^n\right)$ for $j = 0,1$ are the first time derivatives of flux computed in point stencil $\mathcal{S}_{i-\ha}^1;$
                \item[a2.] $f_{i+j}^{(1)} = \frac{1}{\Delta t} \left(f\left(U_{i+j}^n + \Delta tU_{i,j}^{(1)}\right) - f_{i+j}^n\right)$ for $j = 0,1$ are the first time derivatives of flux computed in point stencil $\mathcal{S}_{i+\ha}^1;$
                \item[a3.] $U_{i-1,j}^{(1)} = -\frac{1}{\Delta x}\left(f_{i}^n - f_{i-1}^n\right) $ for $j = 0,1$ are the first time derivatives of the solution $U$ at time $t^n$ for each position of stencil $\mathcal{S}_{i-\ha}^1$ necessary to compute the Taylor expansion truncated at first term;
                \item[a4.] \label{u1coincide} $U_{i,j}^{(1)} = -\frac{1}{\Delta x}\left(f_{i+1}^n - f_{i}^n\right) $ for $j = 0,1$ are the first time derivatives of the solution $U$ at time $t^n$ for each position of stencil $\mathcal{S}_{i+\ha}^1$ necessary to compute the Taylor expansion truncated at first term.
            \end{itemize}
        \end{enumerate}
\end{itemize}
\end{comment}

Finally, the expanded form of the fluxes (\ref{eq:fluxCAT2}) is given by
\begin{align}
\label{F2-}
    F_{i-\ha}^1 &= \frac{1}{4}\left(f^n_{i-1} + f^n_i + f\left(u_{i-1}^n + \Delta t \, u_{i-1,0}^{(1)}\right) + f\left(u_{i}^n + \Delta t \, u_{i-1,1}^{(1)}\right) \right),\\
    \label{F2+}
    F_{i+\ha}^1 &= \frac{1}{4}\left(f^n_i + f^n_{i+1} + f\left(u_{i}^n + \Delta t \, u_{i,0}^{(1)}\right) + f\left(u_{i+1}^n + \Delta t \, u_{i,1}^{(1)}\right) \right),
\end{align}
and the solution is obtained by substituting these fluxes in formula (\ref{CAT2P}).\\ 
The computation of the interfacial flux $F_{i+\ha}^1$ can be recast into the algorithm:
\begin{itemize}
    \item[Step 1:] Compute $f_{i+\ha}^{(0)}$ adopting an interpolation formula over stencil $\mathcal{S}_{i+\ha}^1$ at time $t^n$.
    \item[Step 2:] Compute the first derivatives $u^{(1)}_{i,p}$ in time through the numerical compact Cauchy-Kovalesky procedure, using  $\partial_t u = -\partial_x f $, with data at time $t^n$. % as done in step a3. and a4.;
    \item[Step 3:] Compute the truncated Taylor expansions: $u^{1,n+1}_{i,p} = u^n_{i+p} + \Delta t\, u^{(1)}_{i,p}$ for $p=0$ and $1$.
    \item[Step 4:] Compute the first time derivatives of the flux using the first difference formulas: 
    \[
        f^{(1)}_{i,p} = \Frac{f \left(u^{1,n+1}_{i,p}\right) - f^n_{i+p}}{\Delta t}.
    \]
    \item[Step 5:] Compute $f_{i+\ha}^{(1)}$ through $f^{(1)}_{i,j}$ adopting an interpolation formula on stencil $\mathcal{S}_{i+\ha}^1;$
    \item[Step 6:] Compute $F_{i+\ha}^1$ as a Taylor expansion: $F_{i+\ha}^1 = f^{(0)}_{i+\ha} + \Frac{\Delta t}{2}f^{(1)}_{i+\ha}$ with (\ref{eq:fluxCAT2}).
\end{itemize}   
%Figure \ref{fig:cat2grid} could help to have a graphic idea of the necessary stencil to compute the right flux reconstruction $F_{i+\ha}^1.$
% --- FIG ---

\begin{comment}
In Figure~\ref{fig:cat2grid} we illustrate the associated stencil.
\begin{figure}[!ht]
\begin{center}
    \includegraphics[scale = 0.8]{FIG/CAT_Chap_Theo_fig/cat2grid.pdf}
    \caption{1D local stencil adopted to compute the flux $F_{i+1/2}^1$. The continuous line represents the $x-$direction while the dotted ones the $t-$direction. $U_{i,p}^{1,n+1},$ for $p=0,1,$ correspond to the first-order truncated Taylor expansions in position $x_{i+p},$ which aces the numerical approximation of the time derivatives of the fluxes $f^{(1)}_{i,p}$ at positions $x_{i+p}.$} \ema{To be checked!}
    \label{fig:cat2grid}
\end{center}
\end{figure} 
% --- FIG ---
\end{comment}

\subsubsection{CAT2P}
Generically, the $2P$ CAT scheme follows formulation \eqref{CAT2P} with the interface fluxes given by \eqref{FP}.
The expression of the right numerical flux of order $2P$ is obtained with formula:
\begin{equation}
    \label{FP+}
    F^P_{i+\ha} = \sum_{k=1}^{2P}\frac{\Delta t^{k-1}}{k!}\mathcal{A}^{0,\ha}_{P}\left(f_{i,*}^{(k-1)},\Delta x\right) = \sum_{k=1}^{2P}\frac{\Delta t^{k-1}}{k!}\sum_{j=-P+1}^P \gamma^{0,\ha}_{P,j}f_{i,j}^{(k-1)}, 
\end{equation}
where the high order time derivatives of the flux are computed following and extending the iterative algorithm presented for CAT2 in the previous section (see also \cite{Carrillo-Pares, CPZMR2020, Macca-Pares} for more details):
\begin{enumerate}
    \item Define $f^{(0)}_{i,j} := f(u_{i+j}^n)$ for all $j = -P+1,\ldots,P;$
    \item For every $k = 1,\ldots,2P-1$:
    \begin{enumerate}
        \item Compute the $k$-th derivative of $u$ at time step $t^n$ for each position $x_{i+j}$ with $j = -P+1,\ldots,P$ through the numerical compact version of the Cauchy-Kovalesky identity \eqref{CK}  as:
        $$u_{i,j}^{(k)} = -\mathcal{A}^{1,j}_{P}\left(f_{i,*}^{(k-1)},\Delta x\right); $$
        \item Compute the Taylor expansion of $u$ in time truncated to term $k$ for all positions $x_{i+j}$ with $j = -P+1,\ldots, P$ at time $t^{n+r}$ with $r = -P+1,\ldots,P$ as:
        $$u(x_{i+j},t^{n+r}) \approx u_{i,j}^{k,n+r} = u^n_{i+j} +\sum_{m=1}^k\frac{(r\Delta t)^m}{m!} u^{(m)}_{i,j}; $$
        \item Compute the $k-$th time derivative of flux for each position $x_{i+j}$ with $j = -P+1,\ldots, P$ at time $t^{n}$ as:
        $$f^{(k)}_{i,j} = \mathcal{A}^{k,j}_{P}\left(f_{i,j}^{k,*},\Delta t\right), $$ where $f_{i,j}^{k,*}$ means that we are applying the $\mathcal{A}$ operator in time and in particular we apply the differentiation formula to the set of flux approximations 
        $$ f_{i,j}^{k,n-P+1},\ldots,f_{i,j}^{k,n+P},$$ 
        in which $f_{i,j}^{k,n+r}=f\left(u_{i,j}^{k,n+r}\right)$ for all $j,r =-P+1,\ldots,P$.
    \end{enumerate}
\end{enumerate}

\begin{rem}
    Observe that the computation of the numerical flux  $F_{i+1/2}^{P}$ requires the approximation of $u$ at the nodes of a space-time grid of size $2P \times 2P$, represented by $u_{i,j}^{k,n+r}$, for $-P+1 \leq j,r =\leq P$ (see Figure~\ref{1D_grid}). 
    The approximations of the solution $u$ at successive times $(n-P+1)\Delta t$, \dots, $(n-1)\Delta t$ are different from the ones  already computed in the previous time steps $t^{n-P},\ldots, t^{n-1}$ which are $u_{i+j}^{n-P}$, \ldots, $u_{i+j}^{n-1}$.  
    In other words, the discretization in time is not based on a multi-step method but on a one-step one. In fact, it can be re-interpreted as a Runge-Kutta method whose stages are $u^{n+r}_{i,j}$, $r = -P+1, \dots, P$, see \cite{Carrillo-Pares,CPZMR2020,Macca-Pares}.
\end{rem}
\begin{rem} \label{rem:local}
    These approximations are local. Indeed, suppose that $i_1+j_1 = i_2+j_2 = \ell,$ i.e. $x_{\ell}>0$ belongs to $\mathcal{S}^P_{i_1+1/2}$ and $\mathcal{S}^P_{i_2+1/2}$ with local coordinates $j_1$ and $j_2$ respectively. Then, $f^{(k)}_{i_1,j_1}$ and $f^{(k)}_{i_2,j_2}$ are, in general, two different approximations of $\partial_t^k f(u)(x_{\ell},t^n).$
\end{rem}


\subsubsection{Computational complexity}\label{sec:Comput_Comp}
In this paragraph we estimate the  computational complexity for the CAT$2P$ scheme. 
The details about the complexity of the algorithm can be found in Appendix~\ref{sec:comp_compl}.
As a summary, the operation count per cell per time step for CAT$2P,$ $P>1,$ applied to the scalar case is $3.5(2P)^3 - 1.5(2P)^2 + (2P)$ flop plus $(2P)^3 - 2(2P)^2 + (2P) + 1$ function evaluations. 
For CAT2 this gives $24$ flop and $3$ function evaluations, while for CAT4 we get $204$ flop and $37$ function evaluations.
\begin{comment} 
\begin{table}[!ht]
    \centering
    \numerikNine
    \begin{tabular}{|cc||c|c|c|}
    \hline
    Scheme & $P$ & Cost per flux & Ratio & Ratio \\
    & & {\footnotesize $C(P)=3.5(2P)^3 - 1.5(2P)^2 + (2P)$} & {\footnotesize$C(P+1)/C(P)$} & {\footnotesize CAT2P/CAT2} \\
    \hline
    %1st order & $0$ & 1 & \\
    CAT2 & $1$ & 29   & 1 & 1 \\
    CAT4 & $2$ & 299  & 10.31 & 10.31\\
    CAT6 & $3$ & 1097 &  3.67 & 37.83 \\
    CAT8 & $4$ & 2711 &  2.47 & 93.48 \\
    CAT10 &$5$ & 5429 &  2.00 &187.2 \\
    $\vdots$ & $\vdots$ & $\vdots$ & $\vdots$ & $\vdots$ \\
    CAT2P & $P$ & $C(P)$ & 1.00 & $\infty$ \\ 
    \hline
    \end{tabular}
    \caption{Cost of the CAT schemes in term of number of operations to update one interface flux as a function of $P$. }
    \label{tab:Cost}
\end{table}
\end{comment}
%
% 
%  ACAT : Adaptive limiter
%
%
\subsection{Adaptive limiter for CAT schemes - ACAT schemes} \label{ssec:beyondCAT}
Although the Compact Approximate Taylor (CAT) schemes are linearly stable in the $L^2$-sense under the usual CFL-$1$ condition, they may produce bounded oscillations close to the discontinuity of the solution. 
To avoid these spurious phenomena, an Adaptive \apriori shock-capturing technique has been developed in \cite{CPZMR2020} and called ACAT. There, the order of the method is locally adapted to the smoothness of the numerical solution by means of indicators which check the regularity of the data for each temporal step. More specifically,  once the approximations  of the solution $u$ at time $t^n$ have been computed, the stencil of the data adopted to actually compute the right flux $F_{i+1/2}^P$ are set to belong to 
 $$  \mathcal{S}_p = \{ u_{i-p +1}^n, \dots, u_{i+p}^n\}, \quad p = 1, \dots, P.  $$
 The selected stencil is the one with maximal length among those in which the solution at time $t^n$ is 'smooth'.
 The smoothness is assessed according to some smoothness indicators: $\psi^p_{i+1/2}$, for any $p = 1, \ldots, P$, which are defined as:
\begin{equation}\label{condiconesS0}
\psi^p_{i+1/2} \approx \left\{
\begin{array}{cl}
1 & \mbox { if $u$ is 'smooth' in $\mathcal{S}_p $,}\\
0 & \mbox {otherwise.} 
\end{array}\right.
\end{equation} 
For this strategy one needs to define a robust first-order flux reconstruction, for instance  Rusanov-, HLL- or HLLC-based flux reconstruction \cite{Toro2009,LeVeque2002book}.
Next, one can employ a TVD flux-limiter for the second-order flux reconstruction to be combined with CAT2, such as, for instance minmod, van Alabada or superbee \cite{LeVeque2002book}.

%\subsubsection{ACAT2 method}
\paragraph{ACAT2}
The expression of the ACAT2 numerical method (for $P=1$) based on a flux limiter (see \cite{LeVeque2002book,LeVeque2007book,Toro2009}) is given by
\begin{equation}\label{ACAT2meth}
     u_i^{n+1} = u_i^n + \frac{\Delta t}{\Delta x}\left(F^{*}_{i-1/2} - F^{*}_{i+1/2} \right) ,
\end{equation}
where the fluxes $F^{*}_{i+1/2}$ are blended as   
\begin{eqnarray}\label{ACAT2_flux} 
    {F}^*_{i\pm1/2} & = & \varphi^1_{i\pm1/2} \, {F}^1_{i\pm1/2} +(1-\varphi^1_{i\pm1/2}) \, {F}^{\text{low}}_{i\pm1/2}.
\end{eqnarray}
${F}^1_{i\pm1/2}$ is the CAT2 flux given by \eqref{F2+}-\eqref{F2-}, while $F^{\text{low}}_{i\pm1/2}$ is the first order flux reconstruction, and, $\varphi^1_{i\pm 1/2}$ is a switch computed by a flux limiter which verifies
\begin{equation} \nonumber
    \varphi^1_{i-1/2} \approx \begin{cases}
        1 &\text{if $\{ u_{i-2}^n, \dots, u_{i+1}^n \}$ is 'smooth'}, \\
        0 & \text{otherwise,}
    \end{cases}
    \quad
    \varphi^1_{i+1/2} \approx \begin{cases}
        1 &\text{if $\{ u_{i-1}^n, \dots, u_{i+2}^n \}$ is 'smooth'}, \\
        0 & \text{otherwise.}
    \end{cases}
\end{equation}
For scalar problems, standard flux limiter functions, $\varphi^1(r)$, such as minmod, superbee, van Leer \cite{Sweby,Kemm2010}, may  be used:
\begin{equation}\label{indicador1}
    \varphi_{i+1/2}^1  =  \varphi^1(r_{i+1/2}), 
\end{equation}
where
\begin{equation}\label{wavefuction}
        r_{i+1/2}=
        %\frac{\Delta_{\rm upw}}{\Delta_{\rm loc}} =
         \left\{
        \begin{array}{cl}
            \displaystyle r^L_{i+1/2}= 
            \frac{u^n_i-u^n_{i-1}}{u^n_{i+1}-u^n_{i}} & \mbox {if } a_{i+1/2}  >0, \\
            \displaystyle  r^R_{i+1/2}= 
            \frac{u^n_{i+2}-u^n_{i+1}}{u^n_{i+1}-u^n_{i}} & \mbox {if } a_{i+1/2}  \leq 0,
        \end{array}\right.
        \quad 
        a_{i+1/2} = \begin{cases}  \displaystyle \frac{f(u^n_{i+1}) - f(u^n_i)}{u^n_{i+1} - u^n_i} & \text{ if $|u^n_i - u^n_{i+1}|> \varepsilon$} ,\\
    f'(u^n_i) & \text{otherwise,}
    \end{cases}
\end{equation}
where $a_{i+1/2}$ is an approximation of the wave speed such as Roe's intermediate speed and 
$\varepsilon$ is a small number (but not too small to avoid numerical cancellation). 
    %$$a_{i+1/2} = \begin{cases}  \displaystyle \frac{f(u^n_{i+1}) - f(u^n_i)}{u^n_{i+1} - u^n_i} & \text{ if $|u^n_i - u^n_{i+1}|> tol;$}\\
    %f'(u^n_i) & \text{otherwise.}
    %\end{cases}$$
    %where $tol$ is a given small tolerance.     
An alternative procedure introduced in \cite{Toro2009} avoids the computation of an intermediate speed by defining $\varphi^1_{i+1/2}=\min( \varphi^1(r_{i+1/2}^R), \varphi^1(r_{i+1/2}^L))$.
%\begin{equation}\label{local_indicators}
%        \varphi^1_{i+1/2}=\min( \varphi^1(r_{i+1/2}^R), \varphi^1(r^L_{i+1/2})).
%\end{equation}
This strategy is easily extended to systems by computing the flux limiter component by component.


%\subsubsection{Smoothness indicators}
\paragraph{Smoothness indicators}
The smoothness indicators used in this work  have been introduced in \cite{CPZMR2020}. Nonetheless we briefly recall their construction for the sake of completeness. 
Given a set of the point values $f_j$ of a function $f$ in  the nodes of the stencil $\mathcal{S}^p_{i+1/2}$, $p \geq 2$ we define $\psi_{i+1/2}^p$ as follows. 
First define the lateral left ($^L$) and right ($^R$) weights $w_{i + 1/2}^{p,L/R}$ as
\begin{equation}
        w_{i+1/2}^{p,L}:=\sum_{j=-p+1}^{-1}(f_{i+1+j}-f_{i+j})^2+\varepsilon,\quad
        w_{i + 1/2}^{p,R}:=\sum_{j=1}^{p-1}(f_{i+1+j}-f_{i+j})^2+\varepsilon,
\end{equation}	
where $\varepsilon=10^{-8}$ is a small quantity only used to prevent the weights to vanish. 
Next, using the half harmonic mean,
$w_{i + 1/2}^p = \frac{w_{i + 1/2}^{p,L} w_{i +1/2}^{p,R}}{w_{i + 1/2}^{p,L}+w_{i+1/2}^{p,R}}$,
one defines the high order smoothness indicator over stencil $\mathcal{S}_p$ by 
\begin{equation}\label{indicadores_locales}
        \psi_{i+1/2}^p:=\left( \frac{w_{i + 1/2}^p}{w_{i + 1/2}^p+\tau_{i + 1/2}^p}\right),   
\quad
\text{with}
\quad
\tau_{i + 1/2}^p = \left( (2p - 1)! \sum^{p}_{j=-p+1}\,\gamma^{2p-1, 1/2}_{p,j} \, f^n_{i+j} \right)^2. 
\end{equation}
\begin{comment}
where 
\begin{align}
\tau_{i + 1/2}^P:=& \left( (2p - 1)! \sum^{p}_{j=-p+1}\,\gamma^{2p-1, 1/2}_{p,j} \, f^n_{i+j} \right)^2,
\end{align}
\begin{align}
\tau_{i + 1/2}^P:=& \left( \Delta^{2p-1}_{i-p +1} f\right)^2,
\quad
\text{with}
\quad
 \Delta^{2p-1}_{i-p +1} f =  (2p - 1)! \sum^{p}_{j=-p+1}\,\gamma^{2p-1, 1/2}_{p,j} \, f^n_{i+j}.
\end{align}
$\Delta^{2p-1}_{i-p +1} f$ represents the undivided difference of  $\{f_{i-p+1}, \dots, f_{i+p}\}$.
%\begin{align}
%        \Delta^{2p-1}_{i-p +1} f =&  (2p - 1)! \sum^{p}_{j=-p+1}\,\gamma^{2p-1, 1/2}_{p,j} \, f^n_{i+j}.
%\end{align}    
\end{comment}
These indicators are such that
\begin{equation}\label{condiconesS}
    \psi^p_{i+1/2} \approx \left\{
    \begin{array}{cl}
    1 & \mbox { if $\{f_j\}$ are 'smooth' in $\mathcal{S}^p_i$},\\
    0 & \mbox {otherwise}, 
    \end{array}\right.
\end{equation}
see  \cite{CPZMR2020} for a precise statement of this property and its proof. 
\begin{rem}
Observe that, if data in the stencil $\mathcal{S}^p_i$ are smooth, then
$w_{i + 1/2}^{p,L} = O(\Delta x^2)$, 
$w_{i + 1/2}^{p,R} = O(\Delta x^2)$, and 
$\tau_{i + 1/2}^p = O(\Delta x^{4p})$.
Since we adopt the harmonic mean
\[
    \frac{1}{w_{i + 1/2}^p} = \frac{1}{w_{i + 1/2}^{p,L}} + \frac{1}{w_{i + 1/2}^{p,R}},
\]
then $w_{i + 1/2}^p = O(\Delta x^2)$. 
As such
\[
    \psi_{i+1/2}^p =  \frac{w_{i + 1/2}^p}{w_{i + 1/2}^p+\tau_{i + 1/2}^p} =  \frac{O(\Delta x^2)}{O(\Delta x^2)+ O(\Delta x^{4p})},
\]
so that $\psi_{i+1/2}^p$ is expected to be close to $1$. 
On the other hand, if there is an isolated discontinuity in the stencil then 
$
\tau_{i + 1/2}^p = O(1),
$
therefore one of the lateral weights is $O(1)$ and the other $O(\Delta x^2)$ so that the harmonic mean implies that  $w_{i + 1/2}^p = O(\Delta x^2)$ and thus:
\[
    \psi_{i+1/2}^p =  \frac{w_{i + 1/2}^p}{w_{i + 1/2}^p+\tau_{i + 1/2}^p}  =  \frac{O(\Delta x^2)}{O(\Delta x^2)+ O(1)}, 
\]
so that $\psi_{i+1/2}^p$ is expected to be close to 0. 
\end{rem}

%\subsubsection{ACAT2P methods}
\paragraph{ACAT2P schemes}
Using these ingredients, the final expression of the ACAT2$P$ for $P>1$ scheme is of the form
\begin{equation}
    \label{ACAT2P}
u_i^{n+1} = u_i^n + \frac{\Delta t}{\Delta x} \left({F}^{\mathcal{P}_i}_{i-\ha} - 
      {F}^{\mathcal{P}_i}_{i+\ha} \right),
\end{equation}  
where
\begin{equation}\label{numfluxacat}
    {F}^{\mathcal{P}_i}_{i\pm1/2} = \begin{cases}
    {F}^{*}_{i\pm1/2} & \text{if $\mathcal{P}_i = \emptyset$;}\\
    {F}^{p_{\max}}_{i\pm1/2} & \text{otherwise.}
    \end{cases}
\end{equation}
Here, $\mathcal{P}_i$ is the set of consecutive indices 
    \begin{equation}
    \mathcal{P}_i= \{  p \in \{2, \dots, P \} \ \text{s.t.}\  \psi^p_{i+ 1/2} \approx 1 \},
    \quad \text{and}    \quad
    p_{\max} = \max(\mathcal{P}_i).
    \label{eq:Ai}
\end{equation}
Moreover $F^*_{i+1/2}$ is the ACAT2 second-order numerical flux given by \eqref{ACAT2_flux}, $F^{p_{\max}}_{i+1/2}$ is the ACAT2$p_{\max}$ numerical fluxes defined in \eqref{FP+}. 
If $\mathcal{P}_i \not= \emptyset$, the order of the flux $F^{p_{\max}}_{i+1/2}$ can range from $2P$ to $4$.
Throughout the paper we clip $\psi$ to $1$ if $\psi\geq 0.95$.
\begin{rem}
Notice that in \eqref{eq:Ai} index $p$ starts from $2$ since it is not possible to determine the smoothness of the data in the two-point stencil $\mathcal{S}^1_{i+1/2}$.
\end{rem}


\subsection{Extension to 2D} \label{ssec:CAT_2D}
In this section we focus on the extension of CAT methods to non-linear two-dimensional systems of hyperbolic conservation laws
\begin{equation}
\label{2Dequ}
{u}_t + {f}({u})_x + {g}({u})_y=0.
\end{equation}
The following multi-index notation will be used:
$$
\bi = (i_1, i_2) \in \mathbb{Z} \times \mathbb{Z},
$$
and
$$
\bzero = (0,0),\quad \bone = (1,1), \quad \bhalf= (\ha, \ha), \quad \bef = (1,0), \quad \bes = (0,1).
$$
We consider Cartesian meshes with nodes
$$\mathbf{x}_{\bi} = (i_1 \Delta x, i_2 \Delta y).$$
Using this notation, the general form of the CAT$2P$ method will be as follows:
\begin{equation}
\label{2Dscheme}
u_\bi^{n+1}=u_\bi^n + \frac{\Delta t}{\Delta x}\left[ {F}_{\bi -  \frac{1}{2}\bef}^P-{F}_{\bi +  \frac{1}{2}\bef}^P\right] +
\frac{\Delta t}{\Delta y}\left[ {G}_{\bi -  \frac{1}{2}\bes}^P-{G}_{\bi +  \frac{1}{2}\bes}^P\right],  
\end{equation}
where the numerical fluxes ${F}_{\bi +  \frac{1}{2}\bef}^P$,  ${G}_{\bi +  \frac{1}{2}\bes}^P$ will be computed using the values of the numerical solution $U_{\bi}^n$ in the $P^2$-point stencil centered at $\mathbf{x}_{\bi + \bhalf} = 
((i_1 + \ha)\Delta x, (i_2 + \ha)\Delta y)$ 
$$
S_{\bi+\bhalf}^P = \{ \mathbf{x}_{\bi + \bj}, \quad \bj \in \mathcal{I}_P  \},
$$
where 
$$
\mathcal{I}_P =\{ \bj = (j_1, j_2) \in \Z \times \Z, \quad -P +1 \leq j_k \leq P, \quad k = 1,2 \}.
$$



\subsubsection{2D CAT2}
In order to show the extension of CAT2$P$ procedure let us start with the expression of the CAT2. The numerical fluxes are constructed as follows:
\begin{align}
\label{2Dfluxes}
F_{\bi +  \frac{1}{2}\bef}^1 = & \frac{1}{4}\left( {f}^{1,n+1}_{\bi,\bzero} + {f}^{1,n+1}_{\bi, \bef} + f^{n}_{\bi} +f^{n}_{\bi + \bef}\right), \\
G_{\bi +  \frac{1}{2}\bes}^1 = & \frac{1}{4}\left( {g}^{1,n+1}_{\bi,\bzero} + {g}^{1,n+1}_{\bi, \bes} + g^{n}_{\bi} +g^{n}_{\bi + \bes}\right), 
\end{align}
where 
\begin{align*}
f^{1,n+1}_{\bi, \bj} &= f\left(u_{\bi + \bj}^n +\Delta t \, u^{(1)}_{\bi, \bj}\right), \\
g^{1,n+1}_{\bi, \bj} &= g\left(u_{\bi + \bj}^n +\Delta t \,  u^{(1)}_{\bi, \bj}\right),
\end{align*}
for $\bj= \bzero, \bef$ in the $x$ direction, and $\bj = \bzero, \bes$ in the $y$ direction.
\begin{rem}
Despite what happens for the 1D reconstruction, the first time derivative of $u,$ $u^{(1)}_{\bi, \bj},$ does not coincide in the 2D-grid points. Indeed, observe  that $u^{(1)}_{\bi, \bzero} \neq u^{(1)}_{\bi, \bef}$ and $u^{(1)}_{\bi, \bzero} \neq u^{(1)}_{\bi, \bes}.$ 

Note that, in the 1D case, $u^{(1)}_{i,0 }= u^{(1)}_{i, 1} $.
\end{rem}
Hence, the first time derivatives $u^{(1)}_{\bi,\bj}$ are so defined: 
\begin{align*}
u^{(1)}_{\bi, \bzero} &= -\frac{1}{\Delta x}\left( f_{\bi + \bef}^n - f_{\bi}^n\right) 
-\frac{1}{\Delta y}\left( g_{\bi + \bes}^n - g_{\bi}^n\right),  \\
u^{(1)}_{\bi, \bef} &= -\frac{1}{\Delta x}\left( f_{\bi + \bef}^n - f_{\bi}^n\right)
-\frac{1}{\Delta y}\left( g_{\bi + \bone}^n - g_{\bi + \bef}^n\right), \\ 
u^{(1)}_{\bi, \bes} &= -\frac{1}{\Delta x}\left( f_{\bi + \bone}^n - f_{\bi + \bes}^n\right) 
-\frac{1}{\Delta y}\left( g_{\bi + \bes}^n - g_{\bi}^n\right), 
\end{align*} 
where 
$$ f_{\bi+\bj}^n = f(u_{\bi+\bj}^n), \quad g_{\bi+\bj}^n =g(u_{\bi+\bj}^n), \quad \forall \bj.$$
Finally, the 2D CAT2 method is so defined:
\begin{equation}
\label{2D_CAT2_cons}
u_\bi^{n+1}=u_\bi^n + \frac{\Delta t}{\Delta x}\left[ {F}_{\bi -  \frac{1}{2}\bef}^1-{F}_{\bi +  \frac{1}{2}\bef}^1\right] +
\frac{\Delta t}{\Delta y}\left[ {G}_{\bi -  \frac{1}{2}\bes}^1-{G}_{\bi +  \frac{1}{2}\bes}^1\right],  
\end{equation}

\subsubsection{2D CAT2P}
The high order CAT2$P$ iterative procedure are computed as follows:

\begin{enumerate}

\item  {Define}
$$
f^{(0)}_{\bi,\bj}=f^n_{\bi + \bj}, \quad  g^{(0)}_{\bi,\bj}=g^n_{\bi + \bj}, \quad \bj \in \mathcal{I}_P.
$$

\item  {For $k = 2, \dots, 2P$:}

\begin{enumerate}
\item Compute
\begin{equation*}
 u^{(k-1)}_{\bi,\bj} = - A^{1,j_1}_{P}(f^{(k-2)}_{\bi,(\bcdot, j_2)}, \Delta x) 
 - A^{1,j_2}_{P}(g^{(k-2)}_{\bi,(j_1, \bcdot)}, \Delta y), \quad \bj \in \mathcal{I}_P.
\end{equation*}
\item Compute 
$$
f^{k-1,n+r}_{\bi,\bj} = f \left(  u^n_{\bi + \bj} + 
\sum_{l=1}^{k-1} \frac{(r \Delta t)^l}{l!} u^{(l)}_{\bi,\bj} \right), \quad  \bj\in \mathcal{I}_P,\;\, r = -P+1, \dots, P.
$$
\item Compute
$$
f^{(k-1)}_{\bi,\bj} =   A^{k-1,0}_{P}( f^{k-1, \bcdot}_{\bi,\bj}, \Delta t),\quad  \bj \in \mathcal{I}_P.
$$
\end{enumerate}

\item Compute
\begin{eqnarray}\label{cat22Dx}
F^P_{\bi+\frac{1}{2}\bef}  &= &\sum_{k=1}^{2P} 
\frac{\Delta t^{k-1}}{k!}A^{0, 1/2}_{P}(\tilde{f}_{\bi,(\bcdot, 0)}^{(k-1)}, \Delta x), 
\\\label{cat22Dy}
G^p_{\bi+\frac{1}{2}\bes}  &= &\sum_{k=1}^{2P} 
\frac{\Delta t^{k-1}}{k!}A^{0, 1/2}_{P}(\tilde{g}_{\bi, (0, \bcdot)}^{(k-1)}, \Delta y).
\end{eqnarray}

\end{enumerate}
The notation used for the approximation of the spatial partial derivatives is the following:
\begin{eqnarray*}
 A^{k,q}_{P}(f_{\bi, (\bcdot, j_2)}, \Delta x) & =&  \frac{1}{\Delta x^k} \sum_{l = -P + 1}^P \gamma^{k,q}_{P,l} f_{\bi, (l,j_2)}\\
 A^{k,q}_{P}(g_{\bi, (j_1,\bcdot{} )}, \Delta y) & =&  \frac{1}{\Delta y^k} \sum_{l = -P + 1}^P \gamma^{k,q}_{P,l} g_{\bi, (j_1,l)}
 \end{eqnarray*}
 
 \begin{rem}
 In the last step of the algorithm above the set $\mathcal{I}_P$ can be replaced by its $(2P -1)$-point subset
 $$
 \mathcal{I}^0_P = \{ \bj = (j_1, j_2)\; \text{ \rm{such that} } \;j_1 = 0 \text{ or } j_2 = 0 \}
 $$
 since only the corresponding values of $\tilde{f}^{(k-1)}_{\bi,\bj}$ are  used to compute the numerical fluxes \eqref{cat22Dx} and
 \eqref{cat22Dy}.
 \end{rem}


% Write stuff on how ACAT is not as good as it could be
\subsection{Discussion} \label{ssec:beyondCAT_discusssion}
The ACAT numerical method presented in previous sections is exhaustively described in \cite{CPZMR2020,Macca-Pares} where numerical results are provided.
The extension to systems of conservation laws is done component-by-component, and the extension to multi-dimensions relies on a direction-by-direction splitting.
However we can point out several defects in the design of such an approach.
%\raphael{@Emanuele: please correct, lower or change this list if not appropriate/accurate/wrong...}\\
% A priori vs a posteriori for the anticipation of new extrema
First of all, the impact of the \apriori smoothness indicators is tremendous. Indeed, they must effectively detect the presence of a discontinuity, but also the occurrence of new ones. While the former is achievable with \apriori limiters, the latter is usually more difficult.
Moreover the smoothness indicators can detect discontinuity only if the mesh is fine enough, otherwise it is hard to distinguish between a discontinuity and a smooth region with a large gradient.
%\rred{In my view, the problem is rather that the smoothness indicators capture correctly the discontinuities only if the mesh is fine enough...}\\
% Cost increase and capability decrease when high orders 
Secondly, the smoothness indicator cost drastically increases with $P$. Moreover, it becomes increasingly more difficult to determine the smoothness of a numerical solution when higher and higher order of accuracy is required.
% Tendency to over react
Thirdly, just like any \apriori limiter, the current one has a natural tendency to over-estimate and over-react to possible spurious troubles. As such, the nominal order of accuracy on smooth solution is not always achieved, unless the grid is really fine. \\
These facts have restricted the effective use of ACAT schemes for orders not greater than $2P=6$. 
Also the cost of the current version of ACAT2P constitutes a genuine limitation in 2D. \\
%Even in the $6$-th order case, the cost and efficiency are already genuine limitations in 2D. %\rred{If this sentence is kept, CATMOOD8 should be used in at least one test here...}
In this work we present an alternative way to couple CAT scheme with an \aposteriori limiting technique, called MOOD, and remove
some of the previously described defects.

%=============================================================
%
%  ADD MOOD ONTO CAT --- HERE IS THE NEW STUFF
%
\section{CAT-MOOD} \label{sec:CATMOOD}
The main objective of this work is to combine the \aposteriori shock capturing technique (MOOD) \cite{CDL1} to this family of one-step spatial and temporal reconstructions of high order CAT schemes \cite{CATMOOD_1D}.  

The MOOD algorithm evaluates \aposteriori the solution of the high-order numerical method using a class of criteria that detect a variety of oscillations, even very small ones. It is therefore possible to ensure that the numerical solution preserves some properties of the exact solutions of the PDE system, such as, for example, positivity, monotonicity and increase of the physical entropy, even in complex cases. 
In addition, this procedure allows to reduce and eliminate the numerical oscillations introduced by the high order methods in the presence of shocks or large gradients as is the case for the CAT methods. 

The basic idea of the MOOD procedure is to apply a high-order method over the entire domain for a time step, then check locally, for each cell $i$, the behavior of the solution via admissibility criteria. If the solution computed in cell $i$ at time $t^{n+1}$ is in accordance with the criteria considered, it is kept, otherwise, it is recomputed with a new numerical method of order lower than the previous one. This operation is repeated until acceptability, or, when the (last) first order scheme is used. This last case occurs when the admissibility criteria fail with any previous reconstruction.

Therefore, the object of this work is to design a cascade of CAT methods in which the order is locally adapted according to \aposteriori admissibility criteria thus creating a new family of adaptive CAT methods called CAT-MOOD schemes. 


\subsection{MOOD admissibility criteria} \label{ssec:MOODdetection}
In this work we select 3 different admissibility criteria 
\cite{CDL2,CDL3} which are invoked onto the candidate numerical solution $\left\{ u_i^{n+1} \right\}_{1\leq i \leq I}$:
\begin{enumerate}
    \item 
    \textit{Computer Admissible Detector (CAD)}: This criterion is responsible to detect 
    undefined or un-representable quantities, usually
    not-a-number \texttt{NaN} or infinity quantity due to some division by zero for instance.
    \item 
    \textit{Physical Admissible Detector (PAD)}: The second detector is responsible for ensuring the physical validity of the candidate solution. The detector reacts to every negative pressure $p$ or density $\rho$ in the computational domain, 
    in compliance with (\ref{eq:physical_states}), since otherwise the solution will create non-physical sound speeds, imaginary time steps and so on. 
    The physicality here is assessed from the point of view of a fluid flow, which limits the generality of the criteria as far as predicted pressures are concerned. 
    Said differently, this physical admissibility criteria must be adapted to the model of PDEs which is solved.
    \item 
    \textit{Numerical Admissible Detector (NAD)}: This criterion corresponds to a relaxed variant of the Discrete Maximum Principle (see \cite{Ciarlet,CDL1})
    $$\min_{c\in \mathcal{C}_i^P}(w^n_c) - \delta_i \le w^*_i \le \max_{c\in \mathcal{C}_i^P}(w^n_c) + \delta_i, $$
    where $\mathcal{C}_i^P = \{-P,\ldots,P\}$ is the local centered stencil of order $2P$ and $\delta_i$ is a relaxation term to avoid problems in flat region. $w^*_i$ is the numerical solution obtained with the scheme of order $2P$ and $\delta_i$ is set as:
    \begin{equation}
        \label{eq:delta}
        \delta_i = \max\Big(\varepsilon_1,\varepsilon_2\left[\max_{c\in \mathcal{C}_i^P}w^n_c - \min_{c\in \mathcal{C}_i^P}w^n_c\right]\Big). 
    \end{equation}
    Here $\varepsilon_1$ and $\varepsilon_2$ are small dimensional constant. 
    This criterion is responsible for guaranteeing the essentially non oscillatory (ENO) character of the solution; that is, no large and spurious minima or maxima are introduced locally in the solution\footnote{In the numerical tests of Sec.~\ref{sec:numerics} we compute the relaxed discrete maximum principle only for density $\rho$ and pressure $p$. No limiting, on the velocity components, $u$, $v$, has been considered.}. 
\end{enumerate}    
If a \texttt{NaN} number is detected by CAD in the candidate solution $w_i^*$, then this cell is sent back for recomputation right away. 
Next, if the candidate solution has some positivity issues, then the PAD test is not passed, and, the cell is also invalid.
At last the cell enters the NAD criteria to test for possible numerical oscillation.
If spurious oscillations have contaminated the candidate solution, $w_i^*$, then it will fail NAD.
The candidate solution is then recomputed if at least one of the previous criteria ordered into a chain, see figure~\ref{fig:MOOD_chain_cascade}-left,
has failed.
As a consequence, the MOOD loop drives the code to locally downgrade the order of accuracy by using an auxiliary scheme of lower accuracy.

%\giovanni{[Very nice explanation about what to do if we don't pass NAD, but we do not explain what we do if PAD is not passed, or if CAD is not passed]}
% ---- FIG ----
\begin{figure}[!ht]
    \centering    
    \includegraphics[width=0.7\textwidth]{FIG/chain}
    \includegraphics[width=0.28\textwidth]{FIG/cascade}
    \caption{Left: Detection criteria of the MOOD procedure for a candidate solution $u_i^*$. \textit{Computer Admissible Detector (CAD)}, \textit{Physical Admissible Detector (PAD)}  and \textit{Numerical Admissible Detector (PAD)} ---
    Right: Cascades of CAT schemes used in the MOOD procedure. Starting from the most accurate one, CAT6, downgrading to lower order schemes, and, at last to a $1$st order accurate scheme employed to ensure robustness.}
   \label{fig:MOOD_chain_cascade}
\end{figure}
% ---- FIG ----

\subsection{CAT scheme with MOOD limiting} \label{ssec:CATMOOD} 
In this work we target to reach a maximal $6$th order of accuracy on the part of the domain where a smooth solution is present. 
This can be achieved with CAT6 scheme that we would like to employ as often as possible.
On the contrary, for cells presenting a discontinuous solution we plan to rely on a $1$st order low accurate but robust scheme, for instance using Rusanov or HLL fluxes, that we would employ only when and where necessary. 
In between these two extremes, the CAT4 schemes of $4$th order of accuracy is inserted and tried when CAT6 fails. If CAT4 scheme also fails, then CAT2 scheme ($2$nd order of accuracy) is further tried.
As such we build several cascades of schemes of decreasing orders but possibly increasing robustness, see figure~\ref{fig:MOOD_chain_cascade}-right.
Notice that, apart from CAT6 and the 1st order parachute scheme, the user can decide if the intermediate schemes (CAT4 and CAT2 here) should be included or not. 
In the numerical results in section~\ref{sec:numerics} only the CAT2 scheme is employed to spare computing resources.
This scheme is referred to as CATMOOD6.

  
\subsection{CATMOOD algorithm} \label{ssec:MOODalgo}
Practically, for the time-step $[t^n,t^{n+1}]$ and for each cell $i$, we define a 'mask', $M_i^{n}\in\{-1,0,1\}$, such that
\bea \label{eq:mask}
M_i^n = \begin{cases}
           1 \quad & \text{if} \; u_i^* \; \text{fails at least one criterion},  \\ 
           -1\quad & \text{if} \; \exists  j \in \mathcal{N}_i \; \text{s.t} \; M_j^n=1,\\ 
           0 \quad & \text{otherwise}.
         \end{cases}
\eea
where $\mathcal{N}_i$ is the set of direct neighbor cells of cell $\omega_i$. $M_i^n=-1$ means that cell $i$ is the neighbor of an invalid cell. \\
The algorithm designed for CATMOOD6 scheme is:
\begin{enumerate}
    \item[Init] Let  $\{u_i^n\}_{1\leq c\leq N_c}$ be the numerical solution at $t = t^n$ over the whole domain $\Omega.$
    \item[CAT6] Let $\{u_i^*\}_{1\leq c\leq N_c}$ be the numerical solution at $t = t^{n+1}$ of order 6 obtained from CAT6 scheme. Set $M_i^n=0$.\\
    For all cell $i$, check if $u_i^*$ satisfies all detection criteria. 
    In this case, then $u_i^{n+1} = u_i^*$ and $M_i^n=0$.
    Otherwise $M_i^n=1$ and, if $M_j^n=0$ then set $M_j^n=-1$ for all $j \in \mathcal{N}_i$.
    \item[CAT4] Only for the troubled cell $i$, i.e $M_i^n\not= 0$, recompute $\{u_i^*\}$ the numerical solution at time $t = t^{n+1}$ of order 4 obtained with CAT4 scheme. \\
   Check if $u_i^*$ satisfies all detection criteria.   
    In this case, then $u_i^{n+1} = u_i^*$ and $M_i^n=0$.
    Otherwise $M_i^n=1$ and, if $M_j^n=0$ then set $M_j^n=-1$ for all $j \in \mathcal{N}_i$.
    \item[CAT2] Only for the troubled cell $i$, i.e $M_i^n\not= 0$, recompute $\{u_i^*\}$ the numerical solution at time $t = t^{n+1}$ of order 2 obtained with CAT2 scheme. \\
   Check if $u_i^*$ satisfies all detection criteria.   
    In this case, then $u_i^{n+1} = u_i^*$ and $M_i^n=0$.
    Otherwise $M_i^n=1$ and, if $M_j^n=0$ then set $M_j^n=-1$ for all $j \in \mathcal{N}_i$.
    \item[1st-ord] Only for the remaining troubled cell $i$, i.e $M_i^n\not= 0$, recompute $\{u_i^*\}$ the numerical solution at time $t = t^{n+1}$ of order 1 obtained with a first order scheme, and set $u_i^{n+1} = u_i^*$.
    %\giovanni{I guess if $\exists i|u^*_i $ does not satisfy the criteria then you have to look for a bug... Of course this is a joke, but it could be a useful warning to check even the bullet-proof 1st order scheme}
\end{enumerate}
For efficiency purposes we may remove the CAT4 step in the cascade.

\subsection{Complexity, cost, convergence, implementation}
The unavoidable extra-cost when using a MOOD procedure is the detection of troubled cells with the admissible criteria from section~\ref{ssec:MOODdetection}.
However the detection criteria are particularly inexpensive to compute compared to smoothness indicators or sensors used for instance in ACAT scheme. Usually these admissibility criteria have a negligible cost. \\
The MOOD limiting procedure always converges because the number of cells and schemes are finite. In the worst case scenario the entire solution is computed successively with all CAT2P schemes up to the 1st order accurate solution.
In the best case scenario the solution from the unlimited CAT6 scheme is accepted without any correction.
Any situation in-between is possible, and generally only few cells need to be recomputed. \\
The detection criteria are fundamental, they must be designed to ensure that, if the mesh is fine enough, a smooth solution computed by CAT6 scheme does not produce any troubled cell. 
They must also ensure that, in the vicinity of strong discontinuity, the robust 1st order scheme is regularly employed to avoid spurious oscillations. \\
Concerning the complexity of CATMOOD, it becomes impossible to estimate \textit{a priori} the cost because the limiting adapts to the underlying flow and the computed numerical solution. In the best case scenario mentioned above, the cost of CATMOOD6 is the one of CAT6, in the worst, the cost of CATMOOD6 is the sum of the cost of all schemes in the cascade.
Generally one observes that the amount of troubled cells is of the order of $0-20\%$ of the total number of cells, which makes the MOOD procedure genuinely competitive compared to existing \textit{a priori} limiters, see \cite{CDL2,ADER_MOOD_14} and the numerical section in this paper. 

\subsection{CATMOOD vs LATMOOD}
In the context of the LAT methods proposed by D. Zorío et al. \cite{ZBM2017}, it is noteworthy that the number of operations per time step is significantly reduced compared to CAT$2P$ methods. This reduction is achieved by performing the Taylor expansions only once per point, while CAT$2P$ requires $2P$ expansions per point. However, it is important to acknowledge that CAT$2P$ exhibits superior stability properties compared to LAT.


Consequently, an intriguing question arises: can LATMOOD yield a more efficient approach than CATMOOD? To explore this question, it is important to consider two crucial points: 
\begin{enumerate}
    \item The local computation of numerical fluxes makes CAT more suitable for implementing MOOD. Indeed, unlike CAT schemes, LAT methods compute temporal derivatives in a global manner, which contradicts the local approach of MOOD. Specifically, in CAT methods, once the solution $u^*$ has been computed, if any of the detectors fail, it is sufficient to recompute all local time derivatives of the fluxes with a lower order of accuracy, without causing any issues in the MOOD approach. Contrarily in the case of LAT, since all the approximations are computed in a non local sense (see Remark~\ref{rem:local}) and since LAT does not use a compact stencil, $4P$ computations have to be updated for each bad cell and approximation $u_{i+p}^{(1)},$ $u_{i+p}^{1,n+1},$ etc., making the scheme computationally inefficient.
    \item From our preliminaries tests we notice that the proportion of low order cells substantially increases when LAT is employed instead of CAT. 
\end{enumerate}
